\clearpage\begingroup\NoSourceHead\SourceLeaf{488b}
\makeatletter\tagsleft@true\let\veqno\@@leqno\makeatother
\renewcommand{\thefootnote}{\arabic{footnote})}\setcounter{footnote}{40}
\ArithmeticHead{31}{R}{GAUSS’ PRINCIPIA GENERALIA THEORIAE FIGURAE FLUIDORUM.}
\noindent\hspace*{\parindent}15. Es wirft sich hier die Frage auf, ob \textsc{Gauss} nur das Verdienst zukommt,
zum ersten Mal einen bereits bekannten Satz (die »\textsc{Green}sche Formel«)
in etwas modifizierter Form auf die Umformung der ersten Variation
des Doppelintegrals \textls[100]{angewandt} zu haben, oder ob man ihm auch die \textsc{Green}sche
Formel selbst zuschreiben darf.

Da ist denn zu sagen, dass — soweit ich habe in Erfahrung bringen
können —, die \textsc{Green}sche Formel für die Ebene hier überhaupt zum ersten
Mal vorkommt (als Spezialfall der Formel (41)), wenn man von der \textsc{Ohm}schen
Formel (11) absieht, die eben doch nur als Vorläufer zu betrachten ist, insofern
bei ihr gerade das Charakteristische, der Begriff des Linienintegrals, fehlt.
Dagegen war der analoge Satz, durch den ein Raumintegral in ein Oberflächenintegral
verwandelt wird, bereits in zwei\footnote{In der \textit{Encyklopädie der Mathematischen Wissenschaften}, Bd. II$_1$, S. 477, Fussnote 78) wird in
diesem Zusammenhang auch auf eine Arbeit von \textsc{Lagrange} verwiesen. Dabei handelt es sich aber nur
um Formeln wie
\[
\iiint L\frac{\partial u}{\partial x}\,dx\,dy\,dz
=\iint Lu\,dy\,dz-\iiint u\frac{\partial L}{\partial x}\,dx\,dy\,dz,
\]
bei denen gerade das Charakteristische, die Umformung des Doppelintegrals in ein Oberflächenintegral, fehlt;
daher kommt diese Arbeit hier für uns ebensowenig in Betracht wie die Formel (5) von \textsc{Euler}. Aus demselben
Grund scheidet auch die ebendort, Fussnote 79) zitierte Arbeit von \textsc{Lamé} vom 8. 5. 1829 aus.} vorausgegangenen Arbeiten vorgekommen:
\clearpage\SourceLeaf{489}
\ArithmeticHead{32}{L}{BOLZA, GAUSS UND DIE VARIATIONSRECHNUNG.}
\noindent\textls[100]{Erstens} hatte \textsc{Gauss} selbst schon 1813 in der bereits oben zitierten
Arbeit \textit{Theoria attractionis etc.}\footnote{\textsc{Gauss}, loc. cit. 36), Art. 3—5.} den Satz
\begin{equation*}
\int\frac{\partial W(x,y,z)}{\partial x}\,d\tau
=\int W(x,y,z)\cos(n,x)\,d\sigma\tag{42}
\end{equation*}
und die beiden entsprechenden für $y$ und $z$ für die speziellen Fälle: $W=x$ und
\[
W=F(r),\quad r=\sqrt{(x-a)^2+(y-b)^2+(z-c)^2}
\]
bewiesen; dabei bedeutet $d\tau$ ein Volumenelement, $d\sigma$ ein Oberflächenelement
und $n$ die äussere Flächennormale. \textls[100]{Zweitens} hatte \textsc{Green} in einer vom
29. März 1828 datierten Arbeit\footnote{\textit{An essay on the application of mathematical analysis to the theories of electricity and magnetism},
Nottingham 1828, abgedruckt im Journal für r. u. a. Mathematik, Bd. 39 (1850), 44 (1852), 47 (1854). Der uns
hier interessierende Satz findet sich in Bd. 44, S. 360—362.} denselben Satz (42) und die beiden entsprechenden
mit $W=V\frac{\partial U}{\partial x}$ bewiesen und daraus den nach ihm benannten
Satz über die Umformung des Integrals
\[
\int\left(\frac{\partial U}{\partial x}\frac{\partial V}{\partial x}
+\frac{\partial U}{\partial y}\frac{\partial V}{\partial y}
+\frac{\partial U}{\partial z}\frac{\partial V}{\partial z}\right)d\tau
\]
abgeleitet. Der gewöhnlich als \textsc{Green}sche Formel bezeichnete Satz \textls[100]{in der
Ebene} kommt bei \textsc{Green} gar nicht vor.

Hiernach dürfte die oben aufgeworfene Frage wohl dahin zu beantworten
sein, dass man \textsc{Gauss} nicht nur die Anwendung der \textsc{Green}schen Formel,
sondern auch diese selbst zuschreiben darf.

16. Wir erwähnen gleich hier zwei weitere Autoren, die sich nur wenig
später als \textsc{Gauss} ebenfalls mit der \textsc{Green}schen Formel beschäftigt, und sie
zur Umformung der ersten Variation des allgemeinen Doppelintegrals (1)
benutzt haben, \textsc{Bordoni} und \textsc{Poisson}. — \textsc{Bordoni} hat (1831) in seinem bereits
früher erwähnten \textit{Calcolo sublime}\footnote{\textsc{Bordoni}, loc. cit. 37), Bd. II, Art. 424.} dieselbe Formel (41) wie \textsc{Gauss} abgeleitet,
\clearpage\SourceLeaf{490}
\ArithmeticHead{33}{R}{GAUSS’ PRINCIPIA GENERALIA THEORIAE FIGURAE FLUIDORUM.}
\noindent allerdings mit einem Zeichenfehler (rechts $+B$ statt $-B$). Er schreibt
zunächst in ungenauer Bezeichnung
\begin{equation*}
\begin{aligned}
\iint\frac{\partial A}{\partial x}\,dx\,dy&=\int[A(x_1,y)-A(x_0,y)]\,dy,\\
\iint\frac{\partial B}{\partial y}\,dx\,dy&=\int[B(x,y_1)-B(x,y_0)]\,dx
\end{aligned}\tag{43}
\end{equation*}
und fährt dann fort: Die beiden Integrale
\[
\int Ay'(t)\,dt,\quad\int Bx'(t)\,dt,
\]
erstreckt über die ganze Begrenzungslinie $P$ (auf der die Bogenlänge $t$ gemessen
ist), sind offenbar den beiden einfachen Integralen auf der rechten
Seite von (43) gleich. Über die Umlaufsrichtung wird nichts gesagt, und eine
Folge dieser Nichtbeachtung derselben ist das falsche Vorzeichen in einer der
beiden so erhaltenen Gleichungen. \textsc{Bordoni} erwähnt \textsc{Gauss} nicht, und in der
Tat ist gerade dieses falsche Vorzeichen wohl der beste Beweis dafür, dass
die \textsc{Bordoni}schen Entwicklungen ohne Kenntnis der \textsc{Gauss}schen Arbeit geschrieben
sind, wofür auch sonst alle Anzeichen sprechen.

Im Gegensatz zu \textsc{Gauss} und \textsc{Bordoni} hat \textsc{Poisson}\footnote{\textsc{Poisson}, loc. cit. 30), S. 300.} in seiner Arbeit vom
Jahre 1831 die \textsc{Green}sche Formel in der Ebene in der folgenden Form abgeleitet
\begin{equation*}
\iint\left(\frac{\partial A}{\partial x}+\frac{\partial B}{\partial y}\right)dx\,dy
=\int(A\cos\alpha+B\cos\beta)\,ds,\tag{44}
\end{equation*}
erstreckt über die Begrenzung des Integrationsbereichs des Doppelintegrals in
der $x,y$-Ebene, mit $ds$ als Bogenelement, während $\alpha,\beta$ die Winkel der
äusseren Normale an die Kurve mit der positiven $x$- bezw. $y$-Axe bedeuten.
\textsc{Poissons} Beweis beruht auf den den \textsc{Gauss}schen Formeln (39) und (40) analogen
Gleichungen
\[
\begin{aligned}
dx&=-\cos\beta^0\,ds^0=+\cos\beta'\,ds'\ldots\\
dy&=-\cos\alpha^0\,ds^0=+\cos\alpha'\,ds'\ldots
\end{aligned}
\]

Die Arbeit von \textsc{Poisson} scheint die früheste Stelle zu sein, an der die
\textsc{Green}sche Formel in der Ebene in der uns heute geläufigen Form auftritt.
Auf die Frage, ob \textsc{Poisson} bei Aufstellung der Formel (44) von der \textsc{Gauss}schen
Arbeit beeinflusst war, werden wir weiter unten (Nr. 23) eingehen.
\par\smallskip\noindent{\small X$_2$ Abh. 5.}\hfill{\small 5}\par
\clearpage\SourceLeaf{491}
\ArithmeticHead{34}{L}{BOLZA, GAUSS UND DIE VARIATIONSRECHNUNG.}
\noindent\hspace*{\parindent}17. Nach dieser Abschweifung kehren wir zu dem Ausdruck (37) für
$\delta U$ zurück, in dem nunmehr das erste Doppelintegral durch das Linienintegral
der Gl. (41) zu ersetzen ist. Setzt man für $A,B,C$ ihre Werte ein und verwandelt
das in (37) übrig bleibende Doppelintegral mittels der Gleichung (36)
wieder rückwärts in ein über die Fläche $U$ genommenes Oberflächenintegral,
so ergibt sich schliesslich der folgende Ausdruck\footnote{Vgl. auch unten Nr. 24 und 25, Eingang.} für die \textls[100]{erste Variation
des Flächeninhalts} $U$
\begin{equation*}
\delta U=-\int\delta e\mathbin{.}\cos(5,7)\mathbin{.}dP
+\int\delta e\mathbin{.}\cos(4,5)\mathbin{.}\left(\frac{1}{R}+\frac{1}{R'}\right)dU;\tag{45}
\end{equation*}
darin bedeutet 7 eine Richtung, die senkrecht zum Element $dP$ ist, die Fläche
$U$ berührt und in der Richtung nach der Fläche $U$ zu gerichtet ist, während
$R,R'$ die beiden Hauptkrümmungsradien in dem betrachteten Punkt der Fläche
bedeuten mit einer der heute üblichen entgegengesetzten Vorzeichenbestimmung.

Aus den Gleichungen (34), (35) und (45) ergibt sich nunmehr nach (25)
der Wert von $\delta W$. Dabei tritt eine Vereinfachung ein, wenn man den »Randwinkel«
$(7,8)=i$ einführt und beachtet, dass
\begin{equation*}
\cos(5,7)=\cos(5,8)\cos i.\tag{46}
\end{equation*}

So erhält \textsc{Gauss} das Schlussresultat
\begin{equation*}
\begin{aligned}
\delta W={}&\int dU\mathbin{.}\delta e\mathbin{.}\cos(4,5)\mathbin{.}\left[z+a^2\left(\frac{1}{R}+\frac{1}{R'}\right)\right]\\
&-\int dP\mathbin{.}\delta e\mathbin{.}\cos(5,8)\mathbin{.}(a^2\cos i-a^2+2\beta^2).
\end{aligned}\tag{47}
\end{equation*}

Wir heben daran besonders den invarianten Charakter der Formel hervor,
invariant in Beziehung auf Koordinatentransformation\footnote{$z$ ist hier nicht als Koordinate aufzufassen, vielmehr als senkrechter Abstand über einer horizontalen
Ebene.} sowohl, als auch in
Beziehung auf die Darstellungsform der Flächen $U$ und $T$. Alle in der Formel
vorkommenden Grössen haben eine unmittelbar mit dem vorgelegten
Problem zusammenhängende geometrische Bedeutung. Das wird erreicht durch
die auf den ersten Blick etwas überraschende Darstellung der Doppelintegrale
als Oberflächenintegrale über die Fläche $U$, der Linienintegrale als Linienintegrale
über die Begrenzung $P$ der Fläche $U$, nicht über deren Projektion
auf die $x,y$-Ebene. Der \textsc{Gauss}sche Beweis der Formeln für $\delta U$, $\delta s$ und $\delta W$
\clearpage\SourceLeaf{492}
\ArithmeticHead{35}{R}{GAUSS’ PRINCIPIA GENERALIA THEORIAE FIGURAE FLUIDORUM.}
\noindent ist zwar später nach der Richtung grösserer Einfachheit, Allgemeinheit und
Strenge mehrfach verbessert worden; aber die Resultate selbst sind an Einfachheit
und Eleganz unübertroffen geblieben: sie tragen den Stempel des absolut
vollendeten an sich.
\par\clearpage\endgroup
